\chapter{Configuration reference}
\label{ch:config}

A run is defined by \file{configs/train\_config.yaml}: one top-level key and
four sections. Generate a copy with every default written out explicitly with

\begin{lstlisting}
\end{lstlisting}

This chapter documents every key that file can contain.

\section{How a value is resolved}

Three sources, highest priority first:

\begin{enumerate}
  \item an explicit command-line flag,
  \item the value in the YAML file,
  \item the built-in default.
\end{enumerate}

That ordering is what lets a single committed configuration be swept from the
shell without being edited or copied:

\begin{lstlisting}
poraque-train --config configs/train.yaml --epochs 500
\end{lstlisting}

\begin{pnote}
Unknown keys are \textbf{rejected}, not ignored. A typo such as
\opt{learn\_rate} would otherwise be dropped silently and the run would quietly
use the default --- producing results that do not match the file that appears
to describe them. The error message names the valid keys for that section.
\end{pnote}

The \emph{resolved} configuration --- defaults, file and command-line overrides
merged --- is written beside the results as \file{<json-stem>\_config.yaml}
(by default \file{models/<name>/log/<name>\_config.yaml}). That is the file worth
keeping: it records what actually ran, overrides included.

\section{Every key is optional}
\label{sec:minimal}

Anything omitted takes its default, so a configuration needs to state only
what the run does \emph{differently}. That is what makes a committed file
readable as the description of one experiment rather than a dump of every
knob: of the eighty settings a full configuration carries,
\file{configs/train.yaml} differs in a handful.

\file{configs/train\_complete\_and\_commented.yaml} is the reference: it lists
the settings worth knowing about, each with the reasoning behind it. Read it
there and copy only what you need into your own file. This chapter is the
exhaustive version of the same thing.

Four configurations ship with the source:

\begin{center}
\begin{tabular}{@{}lp{8.4cm}@{}}
\toprule
File & Purpose \\
\midrule
\file{configs/train.yaml} & The clean starting point. Copy this one. \\
\file{configs/train\_complete\_and\_commented.yaml} & The same run, with every
choice explained at length --- the reference to read, not to copy. \\
\file{configs/train\_materialsproject.yaml} & Training on a Materials Project
density download. \\
\file{configs/train\_mixed.yaml} & One operator from several datasets
at once. \\
\bottomrule
\end{tabular}
\end{center}

\section{Top level}

\begin{center}
\begin{tabular}{@{}llp{7.0cm}@{}}
\toprule
Key & Default & Meaning \\
\midrule
\opt{task} & \opt{all} & Which map to train: \opt{ext2chg}, \opt{chg2tau}, or
\opt{all} for both in sequence. \\
\bottomrule
\end{tabular}
\end{center}

\section{\opt{data} --- where the fields come from}

\begin{center}
\begin{tabular}{@{}llp{6.7cm}@{}}
\toprule
Key & Default & Meaning \\
\midrule
\opt{data\_paths} & \file{data/vasp/structures} & The dataset, as a list of
directories --- see below. \\
\opt{cache} & \opt{data/cache} & Where downsampled copies are written. \\
\opt{pattern} & \opt{""} & Prefix filter on the material subdirectories;
empty takes all of them. \\
\opt{format} & \opt{auto} & \opt{vasp}, or \opt{auto} to detect the code from
the files present. Those are the only two values. \\
\opt{resolution} & \opt{32} & Longest grid axis after spectral downsampling. \\
\opt{potcar\_dir} & \opt{null} & \file{POTCAR} library, used where the data
ships no pseudopotentials. \\
\opt{sigma} & \opt{null} & Gaussian pseudo-ion width in \AA{}, where the
Gaussian model is reached. \\
\opt{gaussian\_blur} & \opt{null} & Gaussian blur width in \AA{} applied to the
computed potential. \\
\opt{blur\_method} & \opt{spectral} & \opt{spectral} or \opt{ndimage}. \\
\opt{cache\_in\_memory} & \opt{auto} & Keep decoded fields in RAM between
epochs --- see below. \\
\bottomrule
\end{tabular}
\end{center}

\subsection{\opt{cache\_in\_memory}}
\label{sec:cachememory}

The largest performance setting in the file, and the only one whose effect is
completely invisible from the results.

With it off, \textbf{every epoch reopens, decompresses and re-parses every
field from disk}. That produces exactly the same numbers as caching does ---
and takes about ten times as long. Measured on a V100 with 115 structures at
$32^3$: \opt{\_\_getitem\_\_} was 59\,\% of the training loop, 78.9\,s over 483
calls, while the GPU sat at 2--4\,\% utilisation waiting on \opt{zlib} and
\opt{numpy.array} over a text generator. Turning the cache on took twenty
epochs from \textbf{168.0\,s to 16.3\,s}, with the validation error identical
to five decimals.

It is not a GPU setting. The parse is the same on Apple Silicon and on the
CPU; what an accelerator adds is a fast device left idle by it.

\opt{auto} estimates the decoded size --- grid points $\times$ itemsize
$\times$ channels, over the input and target fields --- and enables the cache
below 4\,GiB, logging what it decided and what it costs:

\begin{lstlisting}
  field cache         : in RAM, ~61.3 MiB (data.cache_in_memory: auto)
\end{lstlisting}

Set it to \opt{false} when the process shares its memory with something the
estimate cannot see. What \opt{auto} refuses on its own is the case the old
default was right about: $128^3$ grids over thousands of materials, where
caching silently would turn a slow run into a dead one.

\begin{pnote}
Parallel loading is \textbf{not} the alternative it looks like. See
\opt{num\_workers} below: added to the cache, workers make training slower.
\end{pnote}

\subsection{\opt{data\_paths}}

\textbf{One key names the dataset, and every entry in it has the same shape:} a
directory of subdirectories, one per material, each holding that material's
volumetric files.

\begin{lstlisting}
data:
  data_paths:
    - data/vasp/structures      # structure_0000/, structure_0001/, ...
    - data/MP                   # mp-124/, mp-81/, ...
    - data/cache/res32          # a cache from an earlier run
\end{lstlisting}

That is what a VASP run tree looks like, what \opt{poraque-mp} writes, and what
the cache builder produces --- so nothing has to be declared about where a path
came from, and no key says so. Every path is pooled into one dataset.

What \emph{differs} is the \textbf{content} of a material's directory, and
content is read, not configured:

\begin{center}
\begin{tabular}{@{}lp{7.4cm}@{}}
\toprule
A material's directory holds & Poraqu\^e \\
\midrule
inputs beside the density & computes $\vext$ from them --- exactly, with a
\file{POTCAR} or \opt{potcar\_dir} \\
a density on its own & computes it from the structure the \file{CHGCAR} carries
in its own header \\
an \file{EXTCAR} already & reads it as it stands (a prepared cache) \\
\bottomrule
\end{tabular}
\end{center}

$\vext$ is \textbf{always computed} for the first two, never read from an
\file{EXTCAR} that happens to sit in a run directory: at inference time there is
no DFT run to read one from, so the input channel has to be Poraqu\^e's own
arithmetic in both places.

\file{TAUCAR} is \textbf{optional, per material}. A directory that has one joins
the \opt{chg2tau} set; one that does not still joins \opt{ext2chg} and is simply
absent from the other. Nothing is declared for it, and a task with no target
anywhere is skipped with a message rather than failing the run --- which is what
makes \opt{type: all} sensible on a Materials Project download, where no $\tau$
exists.

A path holding a \textbf{single} material's run directly --- its files at the
top level rather than one level down --- is read as that one material, so a lone
calculation needs no wrapper directory.

\begin{pnote}
Three keys used to answer this one question: \opt{train\_paths} (a list),
\opt{root} (a single path used when the list was empty) and \opt{source} (which
layout each path was). All three were removed on 2026-08-31. The last is the one
that mattered: it asked the \emph{config} to declare something the
\emph{directory} already answers, which made ``a VASP run'' and ``a Materials
Project download'' two different things to write down when they are the same
thing on disk. A config using any of them now fails and is told what to write
instead.
\end{pnote}

Chapter~\ref{ch:data} covers what each kind of directory contains;
\opt{potcar\_dir} and its Gaussian fallback are in
Section~\ref{sec:mp-potcar-dir}, which is the setting to read before quoting
any number from a run on archive data.

\subsection{\opt{resolution}}

Fields are reduced to this many points along their longest axis before
training. The reduction is a Fourier truncation --- the exact band-limited
projection for a plane-wave field --- so periodicity is preserved and the
electron count is unchanged to machine precision. Interpolation would alias
and shift it.

Cost scales as the cube: raising $32 \to 64$ is roughly eight times the memory
and time. Grid \emph{shapes} are reduced in proportion, so a
$120\times128\times128$ calculation becomes $30\times32\times32$ rather than
being forced square.

\begin{pnote}
There is no key for supplying an external potential. \file{EXTCAR} is
\textbf{always} computed by \pyclass{ExternalPotential} from the \file{POTCAR}
tables, which reproduces a reference potential to a relative error of order
$10^{-5}$. Any \file{EXTCAR} present in a source directory is ignored --- the
training input must be exactly what the pipeline produces at inference time,
when no such file exists. Standard VASP does not write one anyway.
\end{pnote}

\subsection{\opt{gaussian\_blur} and \opt{blur\_method}}

Smooths the computed external potential. The blur is applied \emph{on top of}
the tabulated pseudopotential, never in place of it.

\begin{pwarn}
\opt{ndimage} uses \opt{scipy.ndimage.gaussian\_filter}, which blurs along
\emph{grid} axes and is therefore anisotropic in Cartesian space unless the
cell is orthogonal --- on a face-centred cubic cell the two methods differ by
$2$--$6\,\%$. \opt{spectral} multiplies by $e^{-G^{2}\sigma^{2}/2}$ using the
true reciprocal metric and stays isotropic on any cell. Prefer the default.
\end{pwarn}

Measured at $\sigma = 0.15$~\AA{} and \opt{resolution: 32}, with one structure
held out and everything else fixed, blurring changed the held-out error by
$0.7\,\%$ and the training time by $0.4\,\%$ --- neither meaningful, since the
grid already band-limits the field. It does remove information near the ionic
cores, where the density peaks. Leave it off unless working at higher
resolution.

\begin{ptip}
The cache directory name encodes these choices, for example
\file{res32\_blur0.15spec}, so changing them cannot silently reuse a cache
built with different settings.
\end{ptip}

\section{\opt{model} --- the operator's shape}

\begin{center}
\begin{tabular}{@{}llp{6.7cm}@{}}
\toprule
Key & Default & Meaning \\
\midrule
\opt{width} & \opt{16} & Channel width of the Fourier layers. \\
\opt{modes} & \opt{8} & Retained Fourier modes per axis. \\
\opt{n\_layers} & \opt{4} & Number of Fourier layers. \\
\opt{projection\_channels} & \opt{48} & Hidden width of the output
projection. \\
\opt{activation} & \opt{gelu} & \opt{gelu}, \opt{relu}, \opt{silu} or
\opt{tanh}. \\
\opt{use\_coordinates} & \opt{true} & Append three fractional-coordinate input
channels. \\
\opt{cell\_conditioning} & \opt{true} & Condition every layer on lattice
descriptors (FiLM). \\
\opt{embedding\_dim} & \opt{32} & Width of that cell embedding. \\
\opt{mode\_selection} & \opt{fixed} & \opt{fixed} mode index, or
\opt{physical} at constant $G_{\max}$. \\
\opt{g\_max} & \opt{null} & Cutoff wavevector in \AA$^{-1}$; required by
\opt{mode\_selection: physical}. \\
\opt{pauli\_residual} & \opt{false} & Structural $\tau\ge\tauvW$ for
\opt{chg2tau}. \\
\opt{pauli\_scale} & \opt{null} & Initial Pauli scale in eV/\AA$^{3}$;
\opt{null} fits it from the training split. \\
\opt{learn\_pauli\_scale} & \opt{true} & Optimise that scale alongside the
backbone. \\
\bottomrule
\end{tabular}
\end{center}

\subsection{\opt{width}, \opt{modes}, \opt{n\_layers}}

These three set the capacity. The parameter count is dominated by the spectral
weights --- $4\,w^{2}m^{3}$ complex numbers per layer for width $w$ and
$m$ modes --- so \opt{modes} is by far the expensive knob: doubling it
multiplies the spectral parameters by eight.

\begin{ptip}
\opt{modes} is a \emph{capacity limit}, not a requirement. On a grid too coarse
to supply that many modes, fewer are used automatically, so one configuration
serves materials whose grids differ in size.
\end{ptip}

\subsection{\opt{use\_coordinates} and \opt{cell\_conditioning}}

\opt{use\_coordinates} appends the three fractional coordinates as extra input
channels, giving the network a positional reference the field alone does not
carry.

\opt{cell\_conditioning} feeds lattice descriptors through an embedding of
width \opt{embedding\_dim} and uses it to modulate each layer's activations
(FiLM). Without it the operator sees the field but not the cell it lives in,
and cannot distinguish two materials whose grids agree while their lattice
vectors do not. Keep both on unless deliberately ablating them.

\subsection{\opt{mode\_selection} and \opt{g\_max}}

\opt{fixed} keeps the lowest \opt{modes} indices on every material. Because the
spacing of reciprocal-lattice points depends on the cell size, that means a
\emph{different physical band} for each material.

\opt{physical} instead keeps every mode below a fixed wavevector $G_{\max}$,
retaining $\lfloor G_{\max} L_i / 2\pi \rfloor$ modes along axis $i$, so every
material contributes the same band of physics. Prefer it when cell sizes vary
widely. It requires \opt{g\_max}.

\subsection{\opt{pauli\_residual} and its scale}

For \opt{chg2tau}, the model predicts
\[
  \tau = \tauvW[\rho] + s\,\mathrm{softplus}\bigl(f_\theta(\rho)\bigr),
\]
so the Hoffmann--Ostenhof bound $\tau \ge \tauvW$ holds by construction and the
network learns only the Pauli term $\tauP$. Against a matched baseline it
improved every fold --- $18.3\,\%$ in mean relative $L^{2}$ --- while training
$17.8\,\%$ \emph{faster}, because the network no longer has to re-derive
$|\nabla\rho|^{2}/8\rho$, roughly $31\,\%$ of $\tau$. Recommended.

\opt{pauli\_scale} sets the initial magnitude $s$; \opt{null} fits it from the
training split, which is the sensible default. \opt{learn\_pauli\_scale}
lets the optimiser refine it.

\begin{pwarn}
The bound is a theorem for all-electron densities, whereas \file{CHGCAR} and
\file{TAUCAR} hold pseudo quantities. Check it on new data before enabling the
head --- the training log reports any reference points that violate it.
\end{pwarn}

\section{Numerical precision}

Two independent settings, because they control different costs:

\begin{center}
\begin{tabular}{@{}llp{7.4cm}@{}}
\toprule
Key & Default & Meaning \\
\midrule
\opt{data.precision} & \opt{float64} & How the volumetric fields are
\emph{stored}: \opt{float16}, \opt{float32} or \opt{float64}. A memory
setting --- a $160^3$ field is 16\,MB in double and 8\,MB in single. \\
\opt{model.precision} & \opt{float32} & What the operator \emph{computes} in:
\opt{float32} or \opt{float64}. \\
\bottomrule
\end{tabular}
\end{center}

Keep \opt{data.precision} at \opt{float64} wherever an integral over the cell
is the result --- the electron count, the energy terms are sums of $N^3$
values. Drop to \opt{float32} when the field is on its way into a network that
computes in single precision anyway.

\opt{model.precision: float64} roughly doubles time and memory and exists to
check that a physical result is not a single-precision artefact: an energy
difference of a few meV assembled from $N^3$ voxel sums has no obvious margin
in float32.

\begin{pwarn}
\opt{model.precision: float64} requires \opt{training.device: cpu}. Apple's
Metal backend does not implement double precision at all --- not slowly,
absent --- and the run refuses at startup rather than failing an hour in.
\end{pwarn}

\section{\opt{training} --- how it is fitted}

\begin{center}
\begin{tabular}{@{}llp{6.5cm}@{}}
\toprule
Key & Default & Meaning \\
\midrule
\opt{valid\_fraction} & \opt{0.2} & Fraction of structures held out for
validation; \opt{0} uses every structure. \\
\opt{enable\_kfold} & \opt{false} & Run K-fold cross-validation instead;
ignores \opt{valid\_fraction}. \\
\opt{k\_folds} & \opt{5} & Number of folds, capped at the structure count. \\
\opt{eval\_epoch} & \opt{10} & Evaluate and log every this many epochs. \\
\opt{early\_stopping} & \opt{100} & Stop after this many epochs without
validation improvement; \opt{0} disables. \\
\opt{epochs} & \opt{300} & Passes over the training set. \\
\opt{batch\_size} & \opt{4} & Maximum samples per grid-shape bucket. \\
\opt{learning\_rate} & \opt{0.002} & AdamW step size. \\
\opt{weight\_decay} & \opt{0.0001} & AdamW weight decay. \\
\opt{scheduler} & \opt{cosine} & \opt{cosine} decay, or \opt{null} for a
constant rate. \\
\opt{grad\_clip} & \opt{1.0} & Global gradient-norm clip; \opt{0} disables. \\
\opt{seed} & \opt{0} & Weight initialisation, batch order and fold shuffling. \\
\opt{device} & \opt{auto} & \opt{auto}, \opt{cuda}, \opt{mps} or \opt{cpu}. \\
\opt{strict\_device} & \opt{false} & Abort instead of falling back to the CPU;
see below. \\
\opt{distributed} & \opt{auto} & \opt{auto} or \opt{off}; multi-GPU over
NCCL when the launcher describes a group. \\
\opt{distributed\_} \opt{timeout} & \opt{30.0} & Minutes before a collective
is declared failed. \\
\opt{num\_workers} & \opt{0} & DataLoader worker processes. Read
\opt{data.cache\_in\_memory} first. \\
\opt{pin\_memory} & \opt{auto} & Page-locked staging; \opt{auto} is true on
CUDA and false elsewhere. \\
\opt{tf32} & \opt{true} & TensorFloat-32 matmul and convolution; CUDA, Ampere
and later. \\
\opt{compile} & \opt{false} & \opt{torch.compile} the training loop; CUDA only,
and a measurement switch. \\
\opt{compile\_mode} & \opt{default} & \opt{default}, \opt{reduce-overhead} or
\opt{max-autotune}. \\
\opt{compile\_dynamic} & \opt{true} & One graph over symbolic shapes rather
than one per shape. \\
\opt{loss} & \opt{relative\_l2} & \opt{absolute\_l2}, \opt{relative\_l2},
\opt{absolute\_h1} or \opt{relative\_h1}. \\
\opt{sobolev\_weight} & \opt{0.1} & Gradient-term weight for the two
\opt{h1} losses. \\
\opt{physics} & all \opt{0.0} & Physics-informed loss weights; see
below. \\
\bottomrule
\end{tabular}
\end{center}

\subsection{The validation split}

There is \textbf{one} training protocol and \textbf{one} variation on it.

By default a run trains a single model per task, holding back
\opt{valid\_fraction} of the structures for validation. That is all there is to
it --- no mode to select, no structures to name.

\opt{enable\_kfold} swaps that for K-fold cross-validation: each fold trains a
fresh model on $K-1$ groups and scores it on the group held out. It answers a
different question --- whether the architecture generalises --- and produces
$K$ models rather than one to deploy. \opt{valid\_fraction} is ignored, since
the folds define the splits. Setting \opt{k\_folds} to the number of structures
gives leave-one-out.

\begin{ptip}
Splitting is always at the \emph{structure} level: whole materials move
together. A voxel-level split would place the same crystal on both sides, and
since neighbouring voxels are strongly correlated the score would look
excellent while saying nothing about transfer to a new material.
\end{ptip}

\begin{pwarn}
\opt{valid\_fraction} defaults to \textbf{0.2}, so an ordinary run reports a
genuine held-out score and \opt{early\_stopping} is active. The trade-off is
that the model has then seen only $80\,\%$ of the data. For the final
deployable artefact set \opt{valid\_fraction: 0} --- accepting that its own
metrics become a \textbf{training fit} carrying no generalisation claim, about
four times better than the cross-validated score on the reference dataset ---
and quote a separate \opt{--kfold} run for the number.
\end{pwarn}

\subsection{\opt{eval\_epoch}}

Evaluate and log every this many epochs. Validation is computed \emph{only} on
those epochs, so raising it on a large validation set is a genuine speed-up
rather than only a quieter log. The final epoch always reports, so a run never
ends without a current number.

\begin{lstlisting}
  progress (every 10 epochs):
    train loss: mean PhysicsInformedLoss per batch   |   val rel L2: held-out error, physical units
          epoch     train loss     val rel L2
    -----------------------------------------
         10/300        0.31745        0.34118
         20/300        0.18902        0.21447  *
\end{lstlisting}

\subsection{\opt{early\_stopping}}

Stop after this many epochs without an improvement in the \textbf{validation}
error, and restore the best weights seen:

\begin{lstlisting}
         30/300        0.03173        0.03659  *
         ...
         38/300        0.02249        0.06568
    stopped early at epoch 38: no improvement in 8 epochs
                                 (best 0.03659 at epoch 30)

  trained 38/300 epochs in 12.0 s   loss 0.8599 -> 0.0225
\end{lstlisting}

\begin{pwarn}
It needs a validation split (\opt{valid\_fraction > 0}). With
nothing held out there is only the training loss, which falls monotonically by
construction and so can never signal that training should stop --- asking for
early stopping anyway \textbf{warns} rather than silently doing nothing and
leaving you believing the run was protected.

The shipped default holds a fifth of the structures out
(\opt{valid\_fraction: 0.2}), so early stopping is active out of the box. It
goes inactive only if you set \opt{valid\_fraction: 0} to train on every
structure, and the run says so rather than appearing to be protected.
\end{pwarn}

Patience is counted in \emph{epochs} but checked only on the epochs where
validation is computed, so a value below \opt{eval\_epoch} behaves like
\opt{eval\_epoch}.

The best weights are restored on exit, so the operator you get is the best one
\emph{measured} rather than merely the last one reached --- stopping partway
down a degrading curve would otherwise hand back the degraded model.

The \opt{*} marks an epoch that improved on the best validation score. Only
epochs on which validation was actually measured can do so --- a checkpoint is
written against a measured score, never an assumed one.

\subsection{\opt{batch\_size}}

Capped \emph{per grid-shape bucket}. Materials whose grids differ in shape
cannot be stacked into one tensor, so they are grouped by shape and batches
drawn within a group. A bucket holding a single material yields batches of one
however large this is set. The training log prints the buckets it found.

That cap is a real ceiling, not a formality: \textbf{above the largest bucket's
size the setting stops meaning anything}. In the 115-structure set measured on
a V100 the largest bucket held 70 materials ($\approx$56 after the split), and
\opt{bs=64} and \opt{bs=128} were the same run --- identical peak memory,
identical error, identical time. With the field cache on, throughput responded
to the batch up to 16 and was flat thereafter.

\opt{peak\_vram\_bytes} and \opt{seconds\_per\_epoch} in the metrics JSON are
what a sweep over this should be read from; before they were recorded there,
the only route to either was sampling \opt{nvidia-smi} from outside the
process.

\subsection{\opt{strict\_device}}

\opt{device: auto} prefers CUDA, then Apple Metal, then the CPU. An explicit
backend that is unavailable \textbf{warns and falls back to the CPU}, so a
configuration written on a machine with a GPU still runs on a laptop.

That is right on a workstation and expensive in a queue.
\opt{strict\_device: true} turns the fallback into an error, and
\textbf{every HPC configuration should set it}. The refusal names the probable
cause --- no GPU visible, an empty \opt{CUDA\_VISIBLE\_DEVICES}, a driver older
than the wheel's CUDA runtime, or a build carrying no kernels for this
architecture --- and prints the full device report into the run's log.

That last case is worth stating on its own, because
\opt{torch.cuda.is\_available()} cannot detect it: it answers ``is there a
driver and a device'', not ``can this binary generate code for this device''.
A \opt{+cu130} wheel on a V100 reports available and then aborts at the first
kernel launch, since CUDA 13 dropped Volta. Poraqu\^e checks the compute
capability against the build's own \opt{arch\_list} and refuses before any of
that. See \S\ref{sec:cuda}.

\subsection{\opt{num\_workers} and \opt{pin\_memory}}

\opt{num\_workers} and \opt{data.cache\_in\_memory} (\S\ref{sec:cachememory})
are \textbf{alternatives, not complements}, and the measurement is
unambiguous. Twenty epochs on 115 structures at $32^3$, on a V100:

\begin{center}
\begin{tabular}{@{}lrr@{}}
\toprule
Variant & Training time & Speedup \\
\midrule
As shipped & 168.0\,s & 1.0$\times$ \\
\opt{cache\_in\_memory: true} & \textbf{16.3\,s} & \textbf{10.3$\times$} \\
\opt{num\_workers: 4}, no cache & 120.2\,s & 1.4$\times$ \\
Both & 18.8\,s & 8.9$\times$ \\
\bottomrule
\end{tabular}
\end{center}

The validation error was identical in all four, as it must be. Each worker is a
\emph{process} with a cache of its own, so it re-parses the whole dataset on
its first epoch --- the parse the cache removes is paid $N$ times instead of
once. Raise \opt{num\_workers} only when the data genuinely does not fit in
RAM, which is the case \opt{cache\_in\_memory: false} exists for.

\opt{pin\_memory} is a CUDA transfer optimisation and nothing else: \opt{auto}
is true there and false everywhere, since it does not help on Metal and some
PyTorch versions warn when asked for it.

\subsection{\opt{tf32} and \opt{compile}}

\opt{tf32} keeps float32's range and drops its mantissa to ten bits inside the
tensor cores. It is a real speedup on Ampere and later, and \textbf{exactly
nothing} before it --- a V100 has no TF32 path --- so it is on by default and
costs nothing where it does not apply.

\opt{compile} is a measurement switch rather than a recommendation, and it is
off. The case for it is the kernel profile: at \opt{width=16, modes=8,
n\_layers=3} on a V100, \textbf{44.4\,\% of GPU time is elementwise work}, and
the four costliest kernels are \opt{copy\_}, \opt{Memcpy DtoD}, \opt{add\_} and
\opt{fill\_} --- memory traffic between kernels that each do very little
arithmetic, which is what fusion removes. Convolution is 23.7\,\%, FFT
13.7\,\%, and GEMM only 6.1\,\%.

The case against is that batches are bucketed by grid shape and a real dataset
has many --- 19 distinct shapes in the set measured --- and Inductor
specialises per shape unless \opt{compile\_dynamic} holds across
\opt{rfftn}/\opt{irfftn} with a varying \opt{s=}. Whether it does is the open
question, and the reason this is a flag rather than a default. Compilation is
charged to the first epoch, where \opt{seconds\_per\_epoch} in the metrics JSON
makes the trade visible: a 60-epoch run that pays 90\,s of compilation to save
20\,s is a loss, and a 2000-epoch run is not.

Both are ignored off CUDA. \opt{compile} says so rather than quietly turning on
Inductor's Metal backend, which is a different proposition and untested here.

\subsection{\opt{distributed} and \opt{distributed\_timeout}}
\label{sec:distributed}

\opt{auto} --- the default --- forms a \opt{DistributedDataParallel} group over
NCCL \textbf{when the environment already describes one}, and does nothing
otherwise. A Slurm step with more than one task, or a \opt{torchrun} launch, is
such an environment; a workstation is not.

It cannot turn a one-process run into a four-GPU one. The launcher decides the
topology and this key decides only whether to believe it, so on a cluster the
setting that matters is in the submission script --- see
section~\ref{sec:multigpu}, which carries the whole of it.

\opt{off} refuses a group inside a multi-task allocation, which is how four
independent single-GPU jobs are run from one submission and how a scaling
result is bisected against a single-device baseline.

Three consequences, all visible in the run log. The \textbf{effective batch
size is \opt{batch\_size} times the world size}, so a four-rank run at
\opt{batch\_size: 10} steps on 40 samples and is not the same optimiser as the
one-rank run it is compared against --- halve the learning rate or quarter the
batch size if the comparison is meant to be of the same optimisation.
\textbf{Rank 0 alone writes} the checkpoint, the metrics, the figures and the
PDF, and alone prints. And \textbf{the batches are split, not the samples}:
\opt{ShapeBucketSampler} groups materials by grid shape so no padding ever
reaches the FFT, and a \opt{DataLoader} takes a \opt{sampler} or a
\opt{batch\_sampler} and never both --- so the bucketing runs first,
identically on every rank, and a real \opt{DistributedSampler} partitions the
resulting \emph{list of batches}. Every rank gets a unique, non-overlapping
subset and the same \emph{number} of batches, which is not tidiness: DDP
all-reduces gradients inside each \opt{backward()}, and a rank that runs out of
batches first leaves the others in a collective that never completes.

\opt{distributed\_timeout} is long on purpose. The first collective happens
after every rank has read its prepared cache, and a cold parallel-filesystem
read of a few hundred densities is minutes; a timeout that fires there reports
itself as a NCCL error and sends the reader looking at the network.

\begin{pnote}
There is no \opt{DataParallel} fallback and no Gloo backend. Both would let a
misconfigured run go quietly slower than a single GPU --- the first by
replicating in one process, the second by distributing across CPU cores inside
a GPU allocation. Without CUDA the group is refused with a warning naming the
probable cause, and the run continues on one device.
\end{pnote}

\subsection{\opt{loss} and \opt{sobolev\_weight}}

Four objectives, on two named axes:

\begin{center}
\begin{tabular}{@{}lll@{}}
\toprule
\opt{loss} & normalised per sample & gradients in the objective \\
\midrule
\opt{absolute\_l2} & no  & no  \\
\opt{relative\_l2} & yes & no  \\
\opt{absolute\_h1} & no  & yes \\
\opt{relative\_h1} & yes & yes \\
\bottomrule
\end{tabular}
\end{center}

\textbf{Relative} divides each sample's error by its own target norm, so
materials whose fields differ by orders of magnitude contribute equally and the
value reads directly as a fraction. \textbf{Absolute} does not, so every
\emph{voxel} counts equally and a denser system contributes proportionally more
gradient --- right for a set whose materials are genuinely comparable in
magnitude, and usually wrong for a heterogeneous one, which is why
\opt{relative\_l2} is the default.

\textbf{H1} adds \opt{sobolev\_weight} times the same error on the spatial
gradients. Worth enabling when the derivative matters --- $\tauvW$ depends on
$\nabla\rho$, so gradient noise is amplified in low-density regions. Both
halves are taken in the \emph{same} norm, so the weight is a pure ratio between
values and derivatives rather than also absorbing a change of scale.

\textbf{The validation column follows the objective.} The progress table
reports \opt{val abs L2}, \opt{val rel L2}, \opt{val abs H1} or
\opt{val rel H1} --- whichever is running --- and the number behind it is the
objective's own data term evaluated on the held-out set:

\begin{lstlisting}
    train loss: mean PhysicsInformedLoss per batch   |   val rel H1: held-out error, physical units
    val rel H1 = rel L2 + 0.1 x the relative L2 of the gradient, matching the
    objective; the final per-structure table below reports plain rel L2
          epoch     train loss     val rel H1
\end{lstlisting}

That is not only a label. Early stopping and the checkpoint are decided on this
number, and while it read \opt{val rel L2} regardless a gradient-constrained run
was selecting the model that minimised a functional it was not training on. The
four are not comparable with each other --- which is why the
\emph{per-structure table at the end of a run reports relative $L^2$ whatever
the objective was}. That is the number to quote, and the one that lets two runs
be put side by side.

\begin{pnote}
\opt{loss: sobolev} was renamed on 2026-08-31. It named the gradient term and
left the norm implicit, and offered no unnormalised form at all. It now raises,
naming \opt{relative\_h1} (what it did) and \opt{absolute\_h1} as its
replacements.
\end{pnote}

\subsection{\opt{physics}}

\begin{center}
\begin{tabular}{@{}lp{8.2cm}@{}}
\toprule
Weight & Constraint it penalises the violation of \\
\midrule
\opt{electron\_count\_weight} & $\int\rho\,\mathrm{d}\rr = N$
(\opt{ext2chg}) \\
\opt{positivity\_weight} & Non-negativity of the predicted field \\
\opt{von\_weizsacker\_weight} & $\tau \ge \tauvW$ (\opt{chg2tau}) \\
\opt{euler\_lagrange\_weight} & Orbital-free stationarity residual
(\opt{ext2chg}) \\
\bottomrule
\end{tabular}
\end{center}

All four default to zero, so the objective is the plain supervised baseline
until one is enabled deliberately. Each term is dimensionless, so a single
weight can serve a heterogeneous dataset.

\begin{pwarn}
Introduce them one at a time against a measured baseline, and only once the
data term has converged --- a physics residual evaluated at random
initialisation is noise. A badly scaled constraint degrades accuracy while
looking principled.
\end{pwarn}

Where a constraint can be imposed structurally, prefer that:
\opt{pauli\_residual} enforces $\tau\ge\tauvW$ exactly, whereas
\opt{von\_weizsacker\_weight} only discourages violating it.

\section{\opt{output} --- what is written}

\begin{center}
\begin{tabular}{@{}llp{6.6cm}@{}}
\toprule
Key & Default & Meaning \\
\midrule
\opt{root} & \file{models} & Parent of the run folder. Everything goes in
\file{<root>/<name>/}; \opt{null} disables \emph{all} output. \\
\opt{checkpoint} & \opt{true} & Write \file{<name>.poraque}. \\
\opt{write\_log} & \opt{true} & Write \file{log/}: the log, the metrics JSON
and the resolved configuration. \\
\opt{plot\_figures} & \opt{true} & Render \file{plots/}. \\
\opt{write\_pdf\_report} & \opt{true} & Typeset \file{report/}. \\
\opt{log}, \opt{json} & \opt{null} & Override the two log paths. \opt{null}
derives them from \opt{task.name}; set them only to write outside the run
folder. \\
\opt{plot\_format} & \opt{png} & \opt{png}, \opt{pdf} or \opt{svg}. \\
\opt{dpi} & \opt{200} & Raster resolution for saved figures. \\
\opt{save\_raw\_plot\_data} & \opt{false} & Write the numbers behind each
figure beside it, so a plot can be redrawn without re-running the model. \\
\bottomrule
\end{tabular}
\end{center}

\subsection{\opt{save\_raw\_plot\_data}}

A figure is an argument someone will later want to make differently --- in a
journal's colours, at a journal's size, with two runs on one axis. Re-running
the model to get the numbers back is the expensive way to do that, so with this
on each figure method writes what it drew, beside what it drew:

\begin{lstlisting}
plots/ext2chg_loss_curves.png       plots/ext2chg_loss_curves.csv
plots/ext2chg_parity.png            plots/ext2chg_parity.csv
                                    plots/ext2chg_parity_bin_edges.csv
plots/ext2chg_s000_field_slice.png  plots/ext2chg_s000_field_slice.npz
\end{lstlisting}

The sidecar shares the \textbf{stem} of its figure, so the pairing is a
property of the directory listing rather than of a convention to remember.
Tabular figures get a CSV; the slice figure gets a compressed \file{.npz},
because what it draws are three 2D arrays and a CSV of a $108\times108$ grid is
neither smaller nor easier to read back.

\begin{description}
  \item[The loss curve] is one row per epoch: \opt{epoch}, \opt{train\_loss},
    and the validation column \emph{named for the norm it holds}
    (\opt{val\_rel\_l2}, or \opt{val\_abs\_h1} for an \opt{absolute\_h1} run).
    Epochs on which validation was not measured leave that cell \textbf{empty}
    rather than interpolated: filling them would put points in the file the run
    never measured. There is no physics column --- \opt{train\_loss} is the
    total the optimiser stepped on, and the per-term breakdown is not recorded.
  \item[The parity plot] is one row per \emph{occupied} bin: \opt{split},
    \opt{reference}, \opt{prediction}, \opt{count}, \opt{density}. Long format,
    because that pivots straight back into a \opt{pcolormesh}; occupied bins
    only, because a 200-bin grid is 40\,000 cells of which a few thousand carry
    anything. The bin edges go in their own file, since a log axis makes them
    unrecoverable from the centres.
  \item[The slice figure] stores the three panels as arrays --- \opt{reference},
    \opt{prediction} and \opt{error} --- together with the colour limits the
    figure used. The error is \emph{stored}, not left to be recomputed: a reader
    who reconstructs it and gets something else then knows the file is
    inconsistent, rather than trusting their own subtraction.
\end{description}

It needs \opt{plot\_figures}. The data is written by the figure methods as they
draw, which is what guarantees the file and the image show the same numbers
rather than two independent computations of them.

The resolved configuration is written next to \opt{json} with the suffix
\file{\_config.yaml}. PDF reports are assembled in a temporary directory that
is removed afterwards, so no \file{.tex}, \file{.aux} or \file{.log} files are
left behind.

\section{\opt{symbolic} --- distilling a formula}

Chapter~\ref{ch:symbolic} covers this feature; the keys are listed here for
completeness. It is off by default and needs the optional \opt{symbolic}
install extra.

\begin{center}
\begin{tabular}{@{}p{4.5cm}p{3.1cm}p{7.4cm}@{}}
\toprule
Key & Default & Meaning \\
\midrule
\opt{enable\_symbolic\_} \opt{distillation} & \opt{false} & Search for a
closed-form expression after training. \\
\opt{target} & \opt{model} & \opt{model} distils the operator's predictions;
\opt{reference} fits the DFT data. \\
\opt{features} & \opt{gga} & Variables given to the engine: \opt{gga}
$(\rho, p, q)$, \opt{enhancement} $(p, q)$, or \opt{raw}. \\
\opt{epsilon} & \opt{1e-8} & Vacuum threshold in $e/a_0^3$; clamps every
denominator and drops voxels below it. \\
\opt{unary\_operations} & \opt{[exp, log, sqrt, abs]} & Unary alphabet, passed
to the engine untouched. \\
\opt{binary\_operations} & \opt{['+', '-', '*', /, \textasciicircum]} & Binary
alphabet. \\
\opt{iterations} & \opt{40} & Search iterations --- the main quality/time
knob. \\
\opt{population\_size} & \opt{33} & Individuals per population. \\
\opt{populations} & \opt{15} & Populations run in parallel. \\
\opt{max\_depth} & \opt{10} & Ceiling on expression depth. \\
\opt{max\_size} & \opt{30} & Ceiling on node count. \\
\opt{parsimony} & \opt{0.0032} & Penalty per node; higher gives shorter,
worse-fitting expressions. \\
\opt{n\_samples} & \opt{4000} & Voxels sampled across the dataset. \\
\opt{seed} & \opt{0} & Seeds the sampling and the engine. \\
\opt{template} & \opt{none} & How the target is factorised: \opt{none},
\opt{thomas\_fermi} or \opt{pauli}. \\
\opt{data\_loss} & \opt{mse} & The unpenalised part of the objective;
\opt{mae} is more robust on a density whose tails span orders of magnitude. \\
\bottomrule
\end{tabular}
\end{center}

\subsection*{\opt{symbolic.physics} --- constraints on the search}

\begin{center}
\begin{tabular}{@{}p{4.9cm}p{2.3cm}p{7.8cm}@{}}
\toprule
Key & Default & Meaning \\
\midrule
\opt{enable} & \opt{true} & Penalise violations \emph{inside} the
evolutionary loop rather than filtering the front afterwards. \\
\opt{positivity\_weight} & \opt{100.0} & $F \ge 0$, on the batch. \\
\opt{thomas\_fermi\_weight} & \opt{100.0} & $F(0,0) = 1$: the uniform gas. \\
\opt{von\_weizsacker\_weight} & \opt{100.0} & $F \to 0$ as $p \to \infty$:
one orbital. \\
\opt{p\_infinity} & \opt{1.0e6} & The reduced gradient standing in for
$p \to \infty$ in that probe. \\
\bottomrule
\end{tabular}
\end{center}

\begin{pwarn}[title={Two blocks named \opt{physics}, and they are not the same}]
\opt{training.physics} constrains the \textbf{neural operator}: charge
conservation, positivity of a predicted density, the Euler--Lagrange residual.
Its terms are added to a training loss over voxels.

\opt{symbolic.physics} constrains a \textbf{candidate algebraic expression}
for the Pauli enhancement factor. Its terms are added to a symbolic-regression
fitness over synthetic probe points.

Separate objectives, on separate objects, evaluated at different times by
different engines. Two of the key names appear in both ---
\opt{positivity\_weight} and \opt{von\_weizsacker\_weight} --- and mean
different things in each, which is exactly why they live in separate blocks
rather than sharing a prefix.
\end{pwarn}

\noindent
Filtering the front after a run only measures how few candidates were
physical; by then the populations have spent their budget converging on forms
that must be discarded. With \opt{enable} on, the objective is
\[
  \mathcal{L} = \mathcal{L}_{\rm data}
  + w_{+}\,\overline{\min(F,0)^2}
  + w_{\rm TF}\,|F(0,0) - 1|
  + w_{\rm vW}\,|F(p_\infty, 0)| .
\]
The first term is evaluated on the batch, the other two on probe points --- so
this needs the engine's \emph{full objective}, not an elementwise loss, which
only ever sees \opt{(prediction, target)} pairs and can never evaluate a
candidate where the data does not sit.

\begin{pnote}
The \opt{loss} column of the reported front is then this \emph{constrained}
objective and is not comparable with an unconstrained run's. The $R^2$ and
relative $L^2$ are computed from the expression itself and are unaffected.
\end{pnote}

\section{Worked examples}

\paragraph{Train the deployable models.}
\begin{lstlisting}
task: all
model:    {pauli_residual: true}
training: {valid_fraction: 0, epochs: 200}
\end{lstlisting}

\paragraph{Measure generalisation.}
\begin{lstlisting}
task: all
training: {enable_kfold: true, k_folds: 5}
\end{lstlisting}

\paragraph{A fast smoke test.}
\begin{lstlisting}
data:     {resolution: 20}
model:    {width: 8, modes: 4, n_layers: 2, projection_channels: 16}
training: {epochs: 10}
\end{lstlisting}

\paragraph{Higher fidelity, if the hardware allows.}
\begin{lstlisting}
data:     {resolution: 64}
model:    {width: 32, modes: 12, n_layers: 4}
training: {epochs: 400, batch_size: 2}
\end{lstlisting}
